\documentclass[11pt,a4paper]{article}
\usepackage[margin=25mm]{geometry}
\usepackage{amsmath,amssymb,amsthm,mathtools}
\usepackage[hidelinks]{hyperref}
\newtheorem{proposition}{Proposition}
\newtheorem{corollary}[proposition]{Corollary}
\newcommand{\StateSpace}{\mathsf H}
\newcommand{\OpT}{\mathcal T}
\newcommand{\OpA}{\mathcal A}
\newcommand{\OpB}{\mathcal B_\lambda}
\newcommand{\OpP}{\mathcal P}
\newcommand{\OpQ}{\mathcal Q}
\DeclareMathOperator{\ran}{ran}
\DeclareMathOperator{\Dom}{D}
\DeclareMathOperator{\Rea}{Re}
\title{Well-posedness of a primal PIE and its adjoint-operator dual}
\author{}
\date{}
\begin{document}
\maketitle
\vspace{-12mm}

This note derives a primal--dual equivalence using the notion of
well-posedness in Definition~5 of Jagt and Peet~\cite{jagt}.
It concerns autonomous semigroup well-posedness; it does not, by itself,
establish input/output admissibility or well-posedness of nonlinear feedback.
The argument is an additional derivation, rather than a theorem quoted
from that paper.

\paragraph{Setting and domains.}
Let $\StateSpace$ be a real or complex Hilbert space and let
$\OpT,\OpA\in\mathcal L(\StateSpace)$. Assume that $\OpT$ is injective and has
dense range. Consequently, $\OpT^*$ is also injective and has dense range.
Consider
\begin{equation}\label{eq:pies}
 \frac{d}{dt}\bigl(\OpT v(t)\bigr)=\OpA v(t),
 \qquad
 \frac{d}{dt}\bigl(\OpT^* z(t)\bigr)=\OpA^*z(t).
\end{equation}
For the reconstructed states $x(t)=\OpT v(t)$ and $y(t)=\OpT^*z(t)$,
define
\begin{align}
 G&:=\OpA\OpT^{-1}, &\Dom(G)&=\ran\OpT,\label{eq:G}\\
 G_{\mathrm d}&:=\OpA^*(\OpT^*)^{-1},
 &\Dom(G_{\mathrm d})&=\ran\OpT^*.\label{eq:Gd}
\end{align}
The inverse $\OpT^{-1}$ is defined only on $\ran\OpT$ and is generally
unbounded in the ambient $\StateSpace$ norm. Well-posedness means that the
corresponding operator in \eqref{eq:G} or \eqref{eq:Gd} generates a
$C_0$-semigroup on $\StateSpace$. Classical reconstructed-state solutions start
in the generator domain; arbitrary initial states in $\StateSpace$ have mild
solutions. One should not infer differentiability of the fundamental
state from differentiability of its reconstruction.

\begin{proposition}\label{prop:equivalence}
Under the assumptions above,
\[
 \{\OpT,\OpA\}\text{ is well-posed}
 \quad\Longleftrightarrow\quad
 \{\OpT^*,\OpA^*\}\text{ is well-posed}.
\]
More precisely, if $G$ generates $S(t)$, choose a real
$\lambda\in\rho(G)$ and define
\begin{equation}\label{eq:B}
 \OpB:=\lambda\OpT-\OpA.
\end{equation}
Then $\OpB$ is boundedly invertible and
\begin{equation}\label{eq:similarity}
 G_{\mathrm d}=\OpB^*G^*(\OpB^*)^{-1},
 \qquad
 \Dom(G_{\mathrm d})=\OpB^*\Dom(G^*).
\end{equation}
The dual reconstructed-state semigroup is
\begin{equation}\label{eq:dual-semigroup}
 S_{\mathrm d}(t)=\OpB^*S(t)^*(\OpB^*)^{-1}.
\end{equation}
\end{proposition}

\begin{proof}
Assume that $G$ generates a $C_0$-semigroup on $\StateSpace$.
Its resolvent set contains all sufficiently large positive real numbers,
so choose such a $\lambda$.
Since $\OpT:\StateSpace\to\Dom(G)$ and
$\lambda I-G:\Dom(G)\to\StateSpace$ are both bijections,
\[
 \OpB=(\lambda I-G)\OpT:\StateSpace\to\StateSpace
\]
is bijective. It is bounded because $\OpT$ and $\OpA$ are bounded.
The bounded inverse theorem therefore gives
$\OpB^{-1}\in\mathcal L(\StateSpace)$. Notice that this does not require
$\OpT^{-1}\in\mathcal L(\StateSpace)$.

The primal resolvent has the factorization
\begin{equation}\label{eq:primal-resolvent}
 R(\lambda,G)=\OpT\OpB^{-1}.
\end{equation}
Because $\StateSpace$ is Hilbert, its adjoint semigroup $S(t)^*$ is strongly
continuous and has generator $G^*$; see, for example,
\cite[p.~67]{raymond}. Taking adjoints in
\eqref{eq:primal-resolvent}, with $\lambda$ real, gives
\begin{equation}\label{eq:adjoint-resolvent}
 R(\lambda,G^*)=(\OpB^*)^{-1}\OpT^*.
\end{equation}
In particular,
\begin{equation}\label{eq:domain-equality}
 \Dom(G^*)=(\OpB^*)^{-1}\ran\OpT^*,
 \qquad
 \OpB^*\Dom(G^*)=\ran\OpT^*.
\end{equation}

To identify the action of the dual generator, first note that
\begin{equation}\label{eq:adjoint-intertwining}
 \OpT^*G^*z=\OpA^*z,\qquad z\in\Dom(G^*).
\end{equation}
Indeed, for every $v\in\StateSpace$ and $z\in\Dom(G^*)$,
$\langle G\OpT v,z\rangle=\langle\OpT v,G^*z\rangle$
and $G\OpT v=\OpA v$.
For $z\in\Dom(G^*)$, \eqref{eq:adjoint-intertwining} implies
\[
 \OpB^*z=\OpT^*(\lambda z-G^*z).
\]
Using this identity, injectivity of $\OpT^*$, and the definition of
$G_{\mathrm d}$,
\begin{align*}
 G_{\mathrm d}\OpB^*z
 &=\OpA^*(\OpT^*)^{-1}\OpT^*(\lambda z-G^*z)\\
 &=\lambda\OpA^*z-\OpA^*G^*z\\
 &=(\lambda\OpT^*-\OpA^*)G^*z
 =\OpB^*G^*z.
\end{align*}
Together with \eqref{eq:domain-equality}, this proves the operator
identity \eqref{eq:similarity}, including its domain.
Bounded similarity preserves $C_0$-semigroup generation, yielding
\eqref{eq:dual-semigroup}.

For the reverse implication, apply the same argument to
$(\OpT^*,\OpA^*)$. These operators satisfy the same injectivity and
density assumptions, and their double adjoints are $(\OpT,\OpA)$.
\end{proof}

\paragraph{Why taking adjoints alone is insufficient.}
Directly from the definition of the unbounded adjoint,
\[
 \Dom(G^*)=\{z\in\StateSpace:\OpA^*z\in\ran\OpT^*\},
 \qquad
 G^*z=(\OpT^*)^{-1}\OpA^*z.
\]
Thus $G^*$ and $G_{\mathrm d}$ have reversed multiplication order
and generally different domains. In particular,
$G_{\mathrm d}=G^*$ is not a valid general identity.
The bounded similarity in \eqref{eq:similarity} is what establishes
the desired equivalence in the reconstructed-state formulation.
In fundamental dual coordinates, $z(t)=S(t)^*z_0$ is the mild
evolution, and for $z_0\in\Dom(G^*)$ it satisfies
$\OpT^*\dot z(t)=\OpA^*z(t)$ classically by
\eqref{eq:adjoint-intertwining}. Consistency with reconstruction follows
from $S_{\mathrm d}(t)\OpT^*=\OpT^*S(t)^*$, using
\eqref{eq:adjoint-resolvent} and resolvent--semigroup commutation.

\begin{corollary}[Growth bounds]
The two reconstructed-state semigroups have the same exponential growth
bound. More concretely, if $\|S(t)\|\leq M e^{\omega t}$, then
\[
 \|S_{\mathrm d}(t)\|
 \leq\|\OpB\|\,\|\OpB^{-1}\|\,M e^{\omega t}.
\]
The converse bound follows from
$S(t)=\OpB S_{\mathrm d}(t)^*\OpB^{-1}$.
Consequently, exponential stability is equivalent, although contractivity
in the original norm need not be preserved.
\end{corollary}

\paragraph{Connection to the well-posedness inequalities in the paper.}
For a bounded self-adjoint coercive $\OpP$, suppose the primal certificate
in Corollary~13 of \cite{jagt} is satisfied:
\begin{align}
 \OpT^*\OpP\OpA+\OpA^*\OpP\OpT
 &\preceq2\omega\OpT^*\OpP\OpT,\label{eq:primal-dissipativity}\\
 (\lambda\OpT-\OpA)(\lambda\OpT^*-\OpA^*)
 &\succeq\varepsilon^2 I,\qquad\lambda>\omega.
 \label{eq:primal-surjectivity}
\end{align}
The second inequality makes $\OpB$ surjective. The first implies that
$\OpB v=0$ forces
$(\lambda-\omega)\|\OpT v\|_{\OpP}^2\leq0$, hence $v=0$.
Thus $\OpB$ is boundedly invertible. Moreover, the primal semigroup is
a quasicontraction in the $\OpP$ norm. Its adjoint is a quasicontraction
in the $\OpP^{-1}$ norm. Applying \eqref{eq:dual-semigroup}, define
\begin{equation}\label{eq:dual-storage}
 \OpQ:=\OpB^{-1}\OpP^{-1}(\OpB^*)^{-1}.
\end{equation}
This operator is bounded, self-adjoint and coercive, and the dual
semigroup is a quasicontraction in the $\OpQ$ norm. Equivalently,
\begin{align*}
 \OpT\OpQ\OpA^*+\OpA\OpQ\OpT^*
 &\preceq2\omega\OpT\OpQ\OpT^*,\\
 (\lambda\OpT^*-\OpA^*)(\lambda\OpT-\OpA)
 &=\OpB^*\OpB\succeq\|\OpB^{-1}\|^{-2}I.
\end{align*}
Hence the abstract primal certificate gives an abstract dual certificate.
There is no assertion here that $\OpQ$ belongs to the same finite-degree
polynomial PI parameterization used in a numerical implementation.
Theorem~12 and Corollary~13 in \cite{jagt} are sufficient tests; their
fixed-parameterization feasibility must not be identified with general
well-posedness.

\paragraph{Application to the filtered systems in the IQC construction.}
The relevant filtered systems are
\[
 R_{\mathrm p}=\Theta\begin{bmatrix}P\\I\end{bmatrix},
 \qquad
 R_{\mathrm d}=\mathbf D(\Theta)
                  \begin{bmatrix}P^\top\\I\end{bmatrix},
 \qquad
 \mathbf D(\Theta):=(\widehat J^*\Theta^\top\widehat J)^{-1}.
\]
Here duality is applied to the linear plant and the IQC filter. No adjoint
of the uncertainty is introduced or required. In general,
$R_{\mathrm d}$ is not the adjoint-operator dual of the whole cascade
$R_{\mathrm p}$: taking the dual reverses the order of a cascade,
whereas the IQC construction also inverts and rotates the filter.
Consequently, Proposition~\ref{prop:equivalence} cannot simply be applied
to the entire primal filtered realization to identify the realization
used on the dual side. The cascades require their own argument.

\paragraph{Standing assumptions from the manuscript remain in force.}
The discussion below retains the manuscript's signal assumptions:
$\Theta$ is bounded and causal, $\Theta$ and $\Theta_{22}$ have bounded
causal inverses, and
\[
 \Gamma:=\Theta_{21}P+\Theta_{22}
\]
is bounded and causally boundedly invertible. These assumptions are
available without further verification here. They establish the
zero-initial-state normalization
\[
 \Theta\begin{bmatrix}P\\I\end{bmatrix}
 =\begin{bmatrix}G_0\\I\end{bmatrix}\Gamma,
 \qquad G_0=(\Theta_{11}P+\Theta_{12})\Gamma^{-1}.
\]
They should not be confused with boundedness of a state transition
operator. For classical trajectories with nonzero initial reconstructed
state $x_{f,0}$, linearity instead gives
\[
 \tilde w=\Gamma w+\mathcal O_\Gamma x_{f,0},
\]
where $\mathcal O_\Gamma x_{f,0}$ denotes the free output response.
This notation does not assume that $\mathcal O_\Gamma$ extends to a
bounded map from the reconstructed state space to finite-horizon $L_2$.
The bound on $\Gamma$ estimates only the forced term. With $w=0$ the
free term remains. Conversely, the normalized homogeneous condition
$\tilde w=0$ requires
\[
 w=-\Gamma^{-1}\mathcal O_\Gamma x_{f,0},
\]
provided the free response lies in the domain of the inverse signal map.
Thus the homogeneous systems before and after normalization are not
identified merely by setting both input signals to zero.

The manuscript's zero-initial-state standing assumption applies to its
robust-performance signal statements. A claim of $C_0$-semigroup
generation additionally concerns arbitrary initial reconstructed
states. Imposing zero initial state as well as zero external input would
not establish this claim.

\paragraph{Homogeneous systems with well-posed components.}
From this point, take autonomous well-posedness of the plant and general
PIE filter as given and set the external input to zero. Initial states
remain arbitrary in the reconstructed state spaces. The homogeneous
filtered system is
\[
 \partial_t(\mathcal T_Pv_P(t))=\mathcal A_Pv_P(t),\qquad
 \partial_t(\mathcal T_\Theta v_\Theta(t))
 =\mathcal A_\Theta v_\Theta(t)+\mathcal B_1\mathcal C_Pv_P(t).
\]
The second term is internal state coupling; setting the external input
to zero does not remove it. No input/output well-posedness assumption
on arbitrary initial states is added here; the standing signal
assumptions above remain in force. The question is whether the augmented
homogeneous pair generates a $C_0$-semigroup.

\paragraph{Why autonomous component well-posedness alone is weaker.}
An explicit PIE counterexample uses $H=L_2(0,1)$ and the PI operator
$(Vf)(s)=\int_0^s f(r)\,dr$. Both component pairs $(V,-I)$ are
well-posed: their reconstructed generator is
\[
 G=-V^{-1}=-\partial_s,\qquad
 D(G)=\{x\in H^1(0,1):x(0)=0\},
\]
which generates the right-shift semigroup $S(t)$ with zero inflow.
Choose the plant output $z(t)=-v_P(t)$ and filter input coefficient $I$.
The cascade equations, all with bounded PI coefficients, are
\[
 \partial_t(Vv_P(t))=-v_P(t),\qquad
 \partial_t(Vv_\Theta(t))=-v_\Theta(t)-v_P(t).
\]
Its reconstructed operator is $[G,0;G,G]$ on $D(G)^2$. On smooth
initial states $(x_0,0)$, its solution must satisfy
\[
 x_P(t)=S(t)x_0,\qquad x_\Theta(t)=tS(t)Gx_0.
\]
For any fixed $0<t<1$, choose unit-norm smooth oscillatory $x_0$ supported
inside $(0,1-t)$. Then $\|S(t)Gx_0\|$ can be arbitrarily large.
Consequently this cascade does not generate a $C_0$-semigroup on $H^2$,
despite well-posedness of both autonomous components. The plant's output
$z=Gx_P$ is not an admissible $L_2$ observation. The counterexample
already has zero external input, so restricting attention to homogeneous
systems does not remove this obstruction.

This example can also satisfy the standing boundedness and
invertibility assumptions on the signal operators. Give the plant no
external forcing or feedthrough, so its zero-initial-state map is
$P=0$, while retaining $z=-v_P$ as its output. Realize the filter by
\[
 \partial_t(Vv_\Theta(t))=-v_\Theta(t)+z(t),\qquad
 \begin{bmatrix}\tilde z(t)\\\tilde w(t)\end{bmatrix}
 =\begin{bmatrix}z(t)\\w(t)\end{bmatrix}.
\]
Its internal state is unobservable at the filter output, so its signal
operator is $\Theta=I$, with $\Theta_{22}=I$ and
$\Gamma=\Gamma^{-1}=I$. It is a hard $J$-factor for the identity basis
with weight $J$, for example relative to the uncertainty $\Delta=0$.
Both component homogeneous pairs remain $(V,-I)$, yet the augmented
homogeneous pair is the ill-posed cascade above. No minimal-realization
or arbitrary-initial-state observation assumption in the manuscript
excludes this example. Hence bounds on $\Gamma$ and its inverse alone
cannot close the homogeneous generation argument for these realizations.
Moreover, the associated dual filtered cascade has no state coupling:
the plant's dual output coefficient is $\mathcal B_P^*=0$, and the
dual inverse filter input coefficient is zero because
$\mathcal C_\Theta=0$. Its homogeneous pair consists of two uncoupled
copies of $(V^*,-I)$, each well-posed by
Proposition~\ref{prop:equivalence}. Thus equality of the normalized
signal maps does not establish well-posedness equivalence of the raw
primal and dual filtered realizations.

\paragraph{General PIE filters: the relevant augmented generator.}
The filter in the manuscript is a general PIE; its descriptor must not
be replaced by the identity without an additional assumption. Write
\[
 \partial_t(\mathcal T_\Theta v_\Theta(t))
 =\mathcal A_\Theta v_\Theta(t)+\mathcal B_1 z(t),
 \qquad z(t)=\mathcal C_Pv_P(t).
\]
For a feedforward cascade with the plant, the augmented pair is
\begin{equation}\label{eq:general-cascade-pair}
 \mathcal T_f=\begin{bmatrix}\mathcal T_P&0\\0&\mathcal T_\Theta\end{bmatrix},
 \qquad
 \mathcal A_f=\begin{bmatrix}
  \mathcal A_P&0\\\mathcal B_1\mathcal C_P&\mathcal A_\Theta
 \end{bmatrix}.
\end{equation}
Assume both descriptors are injective with dense range. In reconstructed
coordinates the candidate generator is
\[
 G_f=\mathcal A_f\mathcal T_f^{-1}
 =\begin{bmatrix}
  G_P&0\\\mathcal B_1 C_P^{\rm rec}&G_\Theta
 \end{bmatrix},\qquad
 D(G_f)=D(G_P)\times D(G_\Theta),
\]
where $G_i=\mathcal A_i\mathcal T_i^{-1}$ and
$C_P^{\rm rec}=\mathcal C_P\mathcal T_P^{-1}$.
Even if both diagonal operators generate semigroups, boundedness of the
PIE coefficients alone does not justify bounded perturbation: the
off-diagonal operator need not extend boundedly to the plant state space.

\paragraph{What component well-posedness does establish automatically.}
For all sufficiently large real $\lambda$, component well-posedness
implies that
\[
 E_P(\lambda):=\lambda\mathcal T_P-\mathcal A_P,
 \qquad
 E_\Theta(\lambda):=\lambda\mathcal T_\Theta-\mathcal A_\Theta
\]
are boundedly invertible. The triangular structure therefore gives
\begin{align*}
 E_f(\lambda)&:=\lambda\mathcal T_f-\mathcal A_f
 =\begin{bmatrix}E_P(\lambda)&0\\
 -\mathcal B_1\mathcal C_P&E_\Theta(\lambda)\end{bmatrix},\\
 E_f(\lambda)^{-1}
 &=\begin{bmatrix}
 E_P(\lambda)^{-1}&0\\
 E_\Theta(\lambda)^{-1}\mathcal B_1\mathcal C_P E_P(\lambda)^{-1}
 &E_\Theta(\lambda)^{-1}
 \end{bmatrix}.
\end{align*}
Thus the augmented range condition is automatic for sufficiently large
$\lambda$. In particular,
\[
 E_f(\lambda)E_f(\lambda)^*
 \succeq\|E_f(\lambda)^{-1}\|^{-2}I.
\]
This does not alone prove generation: having a resolvent on a right
half-line is weaker than the semigroup generation estimates.

Consequently, under the assumed component well-posedness, a sufficient
remaining condition is a bounded self-adjoint coercive $\mathcal P_f$
and $\omega\in\mathbb R$ satisfying
\begin{equation}\label{eq:homogeneous-cascade-dissipativity}
 \mathcal T_f^*\mathcal P_f\mathcal A_f
 +\mathcal A_f^*\mathcal P_f\mathcal T_f
 \preceq2\omega\mathcal T_f^*\mathcal P_f\mathcal T_f.
\end{equation}
Choose a sufficiently large $\lambda>\omega$ and apply the generation
criterion in \cite{jagt}. The range condition just proved completes the
argument. This removes the need for an independent augmented range
certificate in this triangular case, but does not make
\eqref{eq:homogeneous-cascade-dissipativity} automatic. In particular,
a semidefinite storage from an indefinite-supply KYP inequality does
not by itself establish this coercive dissipativity estimate.

Another simple sufficient condition is a bounded operator $L$ such that
$\mathcal B_1\mathcal C_P=L\mathcal T_P$. Then the reconstructed
off-diagonal coupling is $L$, and the homogeneous cascade is a bounded
perturbation of $\operatorname{diag}(G_P,G_\Theta)$.

\paragraph{A sufficient cascade condition for general PIE filters.}
Suppose $G_P,G_\Theta$ generate semigroups $S_P(t),S_\Theta(t)$ and
$C_P^{\rm rec}$ is a finite-time $L_2$-admissible observation:
for every $\tau>0$ there is $c_\tau<\infty$ such that
\begin{equation}\label{eq:observation-admissibility}
 \int_0^\tau\|C_P^{\rm rec}S_P(t)x\|^2\,dt
 \leq c_\tau^2\|x\|^2,\qquad x\in D(G_P).
\end{equation}
Then the output trajectory extends continuously to all initial states,
and the cascade semigroup is
\[
 S_f(t)=\begin{bmatrix}S_P(t)&0\\W(t)&S_\Theta(t)\end{bmatrix},
 \qquad
 W(t)x=\int_0^t S_\Theta(t-r)\mathcal B_1
                   C_P^{\rm rec}S_P(r)x\,dr.
\]
Indeed, for $0\leq t\leq\tau$,
\[
 \|W(t)x\|\leq
 \sup_{0\leq r\leq\tau}\|S_\Theta(r)\|
 \|\mathcal B_1\|\sqrt{t}\,c_\tau\|x\|.
\]
This proves boundedness and continuity at zero. Splitting the integral
proves the semigroup identity. Differentiating at zero on
$D(G_P)\times D(G_\Theta)$ identifies its generator there with $G_f$.
To verify equality of domains, choose a sufficiently large real $\lambda$.
The triangular operator $\lambda I-G_f$ is bijective, with inverse
\[
 \begin{bmatrix}
 R(\lambda,G_P)&0\\
 R(\lambda,G_\Theta)\mathcal B_1C_P^{\rm rec}R(\lambda,G_P)
 &R(\lambda,G_\Theta)
 \end{bmatrix}.
\]
Here $C_P^{\rm rec}R(\lambda,G_P)
=\mathcal C_P(\lambda\mathcal T_P-\mathcal A_P)^{-1}$ is bounded.
Since the generator of $S_f$ extends $G_f$ and also has $\lambda$ in its
resolvent, the two operators coincide. This proves state well-posedness
for a genuinely unbounded filter generator. Condition
\eqref{eq:observation-admissibility} is a sufficient condition, not an
automatic consequence of bounded PIE coefficients. Full input/output
well-posedness of the plant includes this observation estimate.

\paragraph{A direct certificate for the augmented PIE.}
Alternatively, apply Corollary~13 of \cite{jagt} directly to
$(\mathcal T_f,\mathcal A_f)$. It suffices to find a bounded self-adjoint
$\mathcal P_f\succeq\alpha I$, with $\alpha>0$, and real numbers
$\omega,\lambda$ and $\varepsilon>0$, with $\lambda>\omega$, such that
\begin{align}
 \mathcal T_f^*\mathcal P_f\mathcal A_f
 +\mathcal A_f^*\mathcal P_f\mathcal T_f
 &\preceq2\omega\mathcal T_f^*\mathcal P_f\mathcal T_f,
 \label{eq:filtered-dissipativity}\\
 (\lambda\mathcal T_f-\mathcal A_f)
 (\lambda\mathcal T_f^*-\mathcal A_f^*)
 &\succeq\varepsilon^2I.\label{eq:filtered-range}
\end{align}
The first inequality establishes dissipativity after shifting by
$\omega$ in an equivalent Hilbert norm. The second gives the range
condition. Together with the descriptor assumptions, these establish
semigroup generation. They are a proposed sufficient verification, not
a certificate already checked for the manuscript's systems. A KYP
inequality with merely semidefinite storage does not automatically
supply these two conditions.

For the general dual inverse filter, disregarding the fixed port
rotations, the descriptor pair is
\[
 \left(\mathcal T_\Theta^*,\,
 \mathcal A_\Theta^*
 -\mathcal C_\Theta^*(\mathcal D_\Theta^*)^{-1}
       \mathcal B_\Theta^*\right).
\]
The inverse correction is bounded at the PIE-coefficient level but its
reconstructed generator contains $(\mathcal T_\Theta^*)^{-1}$.
Consequently, well-posedness of the adjoint-operator filter does not
alone establish well-posedness of this inverse realization. Verify
admissible feedback or apply the generation certificate to this pair,
and then justify its cascade with $P^\top$. One may instead apply
\eqref{eq:filtered-dissipativity}--\eqref{eq:filtered-range} directly to
the full dual filtered pair. For controller interconnections, use the
actual augmented pair, including all controller couplings, rather than
the feedforward pair \eqref{eq:general-cascade-pair}.

\paragraph{Special case only: ordinary and spatially lifted rational filters.}
Write the plant as
\[
 \partial_t(\mathcal T_P v(t))
    =\mathcal A_Pv(t)+\mathcal B_Pw(t),\qquad
 z(t)=\mathcal C_Pv(t)+\mathcal D_Pw(t),
\]
and suppose $G_P=\mathcal A_P\mathcal T_P^{-1}$ on
$D(G_P)=\ran\mathcal T_P$ generates a $C_0$-semigroup.
The assumptions on $\mathcal T_P$ and bounded PIE coefficients are as
in Proposition~\ref{prop:equivalence}. Consider a filter with
$\mathcal T_\Theta=I$:
\[
 \dot\xi(t)=F\xi(t)+B_1z(t)+B_2w(t),\qquad
 \begin{bmatrix}\tilde z(t)\\\tilde w(t)\end{bmatrix}
 =C_\Theta\xi(t)+D_\Theta
                    \begin{bmatrix}z(t)\\w(t)\end{bmatrix},
\]
where $F,B_1,B_2,C_\Theta,D_\Theta$ are bounded on their respective
spaces. This covers rational filters with finitely many states and
their pointwise lifts over the spatial domain. The state space may therefore
be infinite-dimensional even though $F$ is bounded.

In reconstructed coordinates $x(t)=\mathcal T_Pv(t)$, the homogeneous
cascade has operator and domain
\begin{equation}\label{eq:cascade-generator}
 \mathfrak G=
 \begin{bmatrix}G_P&0\\K&F\end{bmatrix},\qquad
 K:=B_1\mathcal C_P\mathcal T_P^{-1},\qquad
 D(\mathfrak G)=D(G_P)\times H_\Theta.
\end{equation}
The coupling $K$ need not be bounded on the ambient plant state space.
Nevertheless, the following argument proves generation without that
extra assumption.

Choose a real $\lambda\in\rho(G_P)$. The operator
$\mathcal B_{P,\lambda}:=\lambda\mathcal T_P-\mathcal A_P$ is
boundedly invertible by the earlier proof. Define
\[
 L:=KR(\lambda,G_P)
   =B_1\mathcal C_P\mathcal B_{P,\lambda}^{-1}
   \in\mathcal L(H_P,H_\Theta).
\]
For $x\in D(G_P)$, this yields $Kx=L(\lambda I-G_P)x$.
The bounded invertible coordinate change
\[
 \mathcal S:=\begin{bmatrix}I&0\\L&I\end{bmatrix},
 \qquad
 \mathcal S^{-1}=\begin{bmatrix}I&0\\-L&I\end{bmatrix}
\]
preserves $D(G_P)\times H_\Theta$ and satisfies
\begin{equation}\label{eq:bounded-cascade}
 \mathcal S\mathfrak G\mathcal S^{-1}
 =\begin{bmatrix}G_P&0\\\lambda L-FL&F\end{bmatrix}
 =\begin{bmatrix}G_P&0\\0&F\end{bmatrix}
  +\begin{bmatrix}0&0\\\lambda L-FL&0\end{bmatrix}.
\end{equation}
The first term generates a $C_0$-semigroup and the second is bounded.
The bounded perturbation theorem~\cite[Thm.~3.4]{schnaubelt} and bounded
similarity therefore show that $\mathfrak G$ generates a $C_0$-semigroup.
This proves autonomous state well-posedness of the primal filtered cascade.
Its input operator
\[
 \begin{bmatrix}\mathcal B_P\\B_1\mathcal D_P+B_2\end{bmatrix}
\]
is also bounded,
so mild state solutions exist uniquely for locally integrable inputs.

For the dual cascade, Proposition~\ref{prop:equivalence} gives the
well-posed plant $P^\top$. If the square filter feedthrough $D_\Theta$
is boundedly invertible, the usual inverse realization of $\Theta^\top$
has bounded state operator
\[
 F_{\mathrm d}=F^*-C_\Theta^*(D_\Theta^*)^{-1}B_\Theta^*,
 \qquad B_\Theta=\begin{bmatrix}B_1&B_2\end{bmatrix}.
\]
The fixed input/output rotations by $\widehat J$ do not change this
state operator. Thus $\mathbf D(\Theta)$ also has a bounded generator,
and the same cascade argument establishes state well-posedness of
$R_{\mathrm d}$. Stability of this inverse filter is a further property;
boundedness of its generator alone proves well-posedness, not stability.
The spatially lifted temporal filter in the numerical examples has
$\mathcal T_\Theta=I$ and bounded $F$, so it falls within the primal
cascade argument above. The inverse-filter argument applies to a square
factor $\Theta$, not to a rectangular basis filter $\Psi$ itself.

\paragraph{General descriptor filters and input/output well-posedness.}
If $\mathcal T_\Theta$ is not boundedly invertible, then
$G_\Theta=\mathcal A_\Theta\mathcal T_\Theta^{-1}$ may be unbounded.
The same coordinate change is not automatically justified: $L$ need not
map into $D(G_\Theta)$, and $G_\Theta L$ need not be bounded.
A simple sufficient condition is that both diagonal operators generate
semigroups and the off-diagonal reconstructed-state coupling extends to
a bounded operator; then bounded perturbation applies directly.
Alternatively, verify the paper's generation inequalities for the full
augmented PIE pair. The inverse dual filter requires its own generation
check when its descriptor is nontrivial.

State generation must also be distinguished from well-posedness of all
input/output maps on the signal spaces used by the IQCs. In particular,
$\mathcal C_P\mathcal T_P^{-1}$ need not be a bounded observation
operator. Its required admissibility, or well-posedness of the complete
plant input/output system, must be established separately. Under those
input/output assumptions the feedforward cascades are causal and
well-defined, since the plant is solved before the filter and no algebraic
feedback loop is introduced by cascading them.

\paragraph{The additional issue when forming $G_0$.}
The normalized graph system
\[
 G_0=(\Theta_{11}P+\Theta_{12})\Gamma^{-1},
 \qquad\Gamma:=\Theta_{21}P+\Theta_{22},
\]
requires inversion of a signal map, rather than just cascading systems.
For a filtered PIE realization
\[
 \partial_t(\mathcal T_f v(t))
    =\mathcal A_fv(t)+\mathcal B_fw(t),\qquad
 \tilde w(t)=\mathcal C_wv(t)+\Lambda w(t),
\]
an invertible $\Lambda$ permits the formal elimination of $w$ and gives
\[
 \mathcal A_0
   =\mathcal A_f-\mathcal B_f\Lambda^{-1}\mathcal C_w.
\]
However, the reconstructed generator perturbation
\[
 -\mathcal B_f\Lambda^{-1}\mathcal C_w\mathcal T_f^{-1}
\]
need not be bounded. Thus invertible feedthrough alone does not
prove generation of $\mathcal A_0\mathcal T_f^{-1}$.
For the zero-initial-state signal argument, the assumed bounded causal
inverse of $\Gamma$ supplies the needed inverse signal map, once the
other constituent maps are well-defined on those signal spaces.
For the particular internal realization, establish admissible feedback,
bounded reconstructed coupling, or a direct generation certificate for
$(\mathcal T_f,\mathcal A_0)$. Transfer-map equality alone does not
prove generation for an arbitrary realization with hidden internal states.
Likewise, a state/input coordinate map involving
$\tilde w=\mathcal C_wv+\Lambda w$ is not automatically a bounded state
similarity on the reconstructed state space.

None of these steps changes the uncertainty in the original feedback
interconnection. Its well-posedness remains a separate hypothesis of the
robust IQC theorem. If a controller is already absorbed into $P$, the
well-posedness assumption on $P$ refers to that controlled nominal
realization; the cascade proof does not independently establish that
controller feedback is admissible.

\begin{thebibliography}{9}
\bibitem{jagt}
D.~S.~Jagt and M.~M.~Peet,
\emph{Verifying Well-Posedness of Linear PDEs using Convex Optimization},
arXiv:2604.00573v2, 2026.
\url{https://arxiv.org/abs/2604.00573v2}.
\bibitem{raymond}
J.-P.~Raymond,
\emph{Evolution Equations: Introduction to Semigroup Theory},
adjoint-semigroup theorem, p.~67.
\url{https://www.math.univ-toulouse.fr/~raymond/evolution-w.pdf}.
\bibitem{shivakumar}
S.~Shivakumar, A.~Das and M.~M.~Peet,
\emph{Dual Representations and $H_\infty$-Optimal Control of
Partial Differential Equations}, arXiv:2208.13104v4, 2026.
\url{https://arxiv.org/abs/2208.13104v4}.
\bibitem{schnaubelt}
R.~Schnaubelt, \emph{Evolution Equations}, lecture notes,
March 19, 2026, Section~3.1.
\url{https://iana.math.kit.edu/downloads/iana3/schnaubelt/Skripten/evgl-skript.pdf}.
\end{thebibliography}
\end{document}
